\scanpage{91}
% Printed seminar programme, rotated 180 degrees relative to the scan.
% A diagonal stroke rules out the entire printed page.
\deleted{SÉMINAIRE NICOLAS BOURBAKI}

\deleted{RÉSUMÉS DES EXPOSÉS}

\deleted{BEAUVILLE Arnaud \quad \textit{Surfaces} K3}

\deleted{La démonstration récente par Y.T. Siu du fait que toute surface K3 est Kählérienne achève une série de travaux qui ont permis de décrire complètement ``l'espace des modules'' des surfaces K3 en termes de leurs périodes. On se propose de décrire les grandes lignes de cette théorie : injectivité de l'application des périodes, surjectivité, cône des classes de Kähler.}

\deleted{CHAPERON Marc \quad \textit{Quelques questions de géométrie symplectique [d'après, entre autres, Poincaré, Arnold, Conley et Zehnder]}}

\deleted{Le dernier théorème géométrique de Poincaré - démontré par G.D. Birkhoff - a conduit V.I. Arnold à formuler de très intéressants problèmes globaux en géométrie symplectique. Dans un travail récent, C.C. Conley et E. Zehnder utilisent une généralisation naturelle de la théorie de Morse pour résoudre l'un d'entre eux, concernant le nombre de points fixes des transformations symplectiques ``homologues à l'identité'' d'une variété symplectique compacte. Une modification presque immédiate de leur méthode permet de répondre à une question voisine sur les intersections de variétés lagrangiennes.}

\deleted{DOLGACHEV Igor \quad \textit{Integral quadratic forms : Algebro-geometrical applications [after V. Nikulin]}}

\deleted{In this talk we will report on the work of V. Nikulin on arithmetics of integral quadratic non-unimodular forms and their applications to surface singularities, automorphisms of K3-surfaces and topology of real algebraic surfaces.}

\deleted{DUFLO Michel \quad \textit{Analyse harmonique sur les groupes algébriques complexes : Formule de Plancherel [d'après M. Andler] et conjecture de M. Vergne}}

\deleted{La méthode des orbites, inaugurée par A.A. Kirillov dans son étude des groupes de Lie nilpotents, permet de décrire un certain nombre des objets traditionnels de l'analyse harmonique sur un groupe de Lie en termes de la représentation co-adjointe. Pour les groupes algébriques complexes, on connait ainsi la formule de Plancherel, et on peut résoudre une conjecture de M. Vergne sur la formule de ``Poisson-Plancherel''.}

\deleted{HAKIM Monique \quad \textit{Valeurs au bord de fonctions holomorphes bornées en plusieurs variables complexes}}

\deleted{Il existe des fonctions holomorphes bornées non constantes dans la boule unité de $\mathbb{C}^{n}$, dont les limites radiales sont presque partout de module 1. On expose ce résultat dû à A.B. Aleksandrov et on étudie des constructions de fonctions holomorphes bornées ayant des propriétés remarquables, introduites par divers auteurs.}

\deleted{SADOURNY Robert \quad \textit{Quelques aspects mathématiques de la dynamique des fluides géophysiques}}

\deleted{La dynamique des fluides géophysiques se partage entre l'expérimentation numérique lourde et quelques approches plus théoriques : existence d'attracteurs dans les espaces de solutions, permettant l'approche de certains mécanismes au moyen d'équations filtrées ; étude, par les méthodes des systèmes dynamiques, des changements de régime de la circulation atmosphérique ou du climat.}
% archive sheet number: 90
